% ============================================================================
% CONTENIDO DESCARTADO / BORRADOR ANTERIOR - Sección LTI
% Extraído de latex.md para limpieza del informe principal.
% Este contenido fue reemplazado por la sección "Linealización por
% Retroalimentación NL: Modelo LTI equivalente" en el informe final.
% ============================================================================

% ------------ LTI (versión anterior, basada en linealización Jacobiana) -----

\subsection{Análisis del Modelo Lineal Invariante en el Tiempo (LTI)}
Si bien el modelo LPV (\textit{Linear Parameter Varying}) presentado anteriormente describe la dinámica global del accionamiento, el diseño y sintonización de controladores lineales clásicos requiere validar el comportamiento del sistema mediante modelos LTI locales. Estos modelos representan la dinámica incremental alrededor de puntos de equilibrio específicos ($\mathbf{x}_{0}$), obtenidos mediante la técnica de linealización Jacobiana.

A continuación, se detalla la deducción analítica de los términos de la matriz del sistema y se evalúa la estabilidad a lazo abierto.

\subsubsection{Linealización del Torque Gravitacional}
La principal fuente de no linealidad en el subsistema mecánico es el torque gravitacional reflejado al eje del motor, $f_g(\theta_m)$. Para incorporar este efecto en la matriz de estado $\mathbf{A}$, se aplica la expansión en Serie de Taylor de primer orden alrededor de un ángulo de operación $\theta_{0}$:

\begin{equation}
    f_g(\theta_m) = -\frac{g \cdot k_l}{J_{eq} \cdot r} \sin\left(\frac{\theta_m}{r}\right) \approx f_g(\theta_{0}) + \underbrace{\left. \frac{\partial f_g}{\partial \theta_m} \right|_{\theta_{0}}}_{A_{2,1}} (\theta_m - \theta_{0})
\end{equation}

El término derivativo define el coeficiente de \textbf{rigidez dinámica} ($A_{2,1}$) en la matriz de estado:
\begin{equation} \label{eq:stiffness}
    A_{2,1}(\theta_{0}) = \frac{\partial f_g}{\partial \theta_m} = -\frac{g \cdot k_l}{J_{eq} \cdot r^2} \cos\left(\frac{\theta_{0}}{r}\right)
\end{equation}
La Ecuación \eqref{eq:stiffness} demuestra que la estabilidad natural del brazo depende de su posición geométrica, variando proporcionalmente al coseno del ángulo de carga.

\subsubsection{Matriz de Estado Jacobiana Completa}
Evaluando las derivadas parciales para el vector de estados completo $\mathbf{x} = [\theta_m, \omega_m, i_{qs}^r, i_{ds}^r, i_{0s}, T_s]^T$, y considerando un punto de equilibrio mecánico estático ($\omega_{m0}=0$), se obtiene la estructura simbólica de la matriz del sistema $\mathbf{A} \in \mathbb{R}^{6 \times 6}$.

\begingroup
\renewcommand{\arraystretch}{1.8}
\begin{equation} \footnotesize
\mathbf{A} =
\begin{bmatrix}
0 & 1 & 0 & 0 & 0 & 0 \\
A_{2,1} & -\frac{b_{eq}}{J_{eq}} & \frac{3 P_p \lambda_r}{2 J_{eq}} & 0 & 0 & 0 \\
0 & -\frac{P_p \lambda_r}{L_q} & -\frac{R_{s0}}{L_q} & 0 & 0 & \Psi_{qT} \\
0 & 0 & 0 & -\frac{R_{s0}}{L_d} & 0 & \Psi_{dT} \\
0 & 0 & 0 & 0 & -\frac{R_{s0}}{L_{ls}} & 0 \\
0 & 0 & \frac{3 R_{s0} i_{q0}}{C_{ts}} & \frac{3 R_{s0} i_{d0}}{C_{ts}} & 0 & \Gamma_{th}
\end{bmatrix}
\end{equation}
\endgroup

Donde $R_{s0}$ es la resistencia a la temperatura de equilibrio. Los términos de acoplamiento térmico-eléctrico se definen como $\Psi_{xT} = -\frac{i_{x0} R_{ref} \alpha_{cu}}{L_x}$ (para $x=q,d$).

El autovalor del subsistema térmico ($\Gamma_{th}$) determina la estabilidad térmica:
\begin{equation}
    \Gamma_{th} = \frac{1}{C_{ts}} \left[ \frac{3}{2} R_{s,ref} \alpha_{cu} (i_{q0}^2 + i_{d0}^2) - \frac{1}{R_{th}} \right]
\end{equation}
Dado que el término de disipación pasiva ($1/R_{th}$) domina sobre la generación incremental de calor para el rango operativo, se verifica que $\Gamma_{th} < 0$, garantizando la ausencia de embalamiento térmico (\textit{thermal runaway}).

\subsubsection{Evaluación en Puntos de Operación Críticos}
Para caracterizar la respuesta dinámica, se calculan los autovalores del subsistema mecánico (submatriz $2\times2$ superior izquierda) en dos configuraciones extremas:

\vspace{0.3cm}

\textbf{Caso A: Equilibrio Estable ($\theta_l = 0^\circ$)} \\
Corresponde al brazo suspendido verticalmente. Al evaluar \eqref{eq:stiffness} con $\cos(0)=1$, el término de rigidez es máximo y negativo (fuerza restauradora).
\begin{itemize}
    \item \textbf{Autovalores:} $\lambda_{1,2} = -0.3345 \pm j 3.5865$
    \item \textbf{Análisis:} La parte real negativa confirma estabilidad asintótica. La parte imaginaria indica una respuesta \textbf{subamortiguada} (comportamiento de péndulo).
\end{itemize}

\vspace{0.2cm}

\textbf{Caso B: Carga Máxima ($\theta_l = 90^\circ$)} \\
Corresponde al brazo horizontal. Aquí $\cos(\pi/2)=0$, anulando el término $A_{2,1}$.
\begin{itemize}
    \item \textbf{Autovalores:} $\lambda_{1} = 0 \quad ; \quad \lambda_{2} = -0.6689$
    \item \textbf{Análisis:} La aparición de un polo en el origen indica la pérdida de rigidez dinámica. El sistema se comporta localmente como un \textbf{integrador puro} (masa inercial con fricción), situándose en el límite de estabilidad marginal.
\end{itemize}

\subsubsection{Conclusión y Modelo de Diseño}
El análisis revela una variación drástica en la dinámica: el sistema pasa de ser un oscilador ($0^\circ$) a un integrador ($90^\circ$). Diseñar un controlador lineal fijo para una planta con polos móviles resulta subóptimo.

Por tanto, para la etapa de diseño se adoptará una estrategia de \textbf{Linealización por Feedback (Feedforward)}. Al compensar el término gravitacional $g(\theta)$ mediante software, se fuerza al sistema a comportarse según el \textit{Modelo de Diseño Nominal} (inercia pura invariante):
\begin{equation}
    G_{mec}(s) = \frac{1}{J_{eq}s + b_{eq}}
\end{equation}

% ------------ Ley de Control Mínima (versión anterior) ----------------------

\subsection{Restricción o Ley de Control Mínima}
Se pretende encontrar la acción mínima para desacoplar el flujo del torque electromagnético, para ello es necesario que la proyección del flujo sobre el eje directo sea nula.
\begin{equation*}
    T_e(t) = \frac{3}{2}P_p(\phi_di_{qs}^r-\phi_qi_{ds}^r)
\end{equation*}
Ya que \[\phi_d = \phi_{PM} +L_di_d, \ \phi_q = L_qi_q\]
Basta con cumplir la condición de \(i_{ds} = 0\) con tal de lograr dicho objetivo, cuando imponemos dicha condición podemos, desarrollar la ecuación que rige a esta corriente.
\[i_{ds} \equiv 0 \implies \dot i_{ds} = 0 \]
\begin{equation}\label{eq:v_ds minima}
    v_{ds}^r= - L_qi_{qs}^r\omega_r(t)
\end{equation}
Como las únicas entradas manipulables en nuestro sistema son los voltajes \(v_a, v_b,v_c\) , es necesario, para lograr la suposición de \(i_{ds}^r=0\), imponer en la entrada \(v_{ds}^r\) la ecuación \ref{eq:v_ds minima} mediante la Transformación Inversa de Park.

\subsubsection{Modelo de la dinámica \textit{residual}}
Si no se cumple la condición inicial esperada para \(i_{ds}^r \) (\(i_{ds}^r(t_0) = 0\) es de esperarse que esta presente cierta dinámica "\textit{residual}", para analizarla, simplemente tomamos de \eqref{eq:equaciones de estado} la dinámica de la corriente en eje directo, tal que:
\begin{equation}\label{eq:corriente id}
    \frac{di_{ds}^r}{dt} = \frac{1}{L_d}\Big[v_{ds}^r - R_s(T_s^\circ(t))i_{ds}^r+L_qi_{ds}^r\omega_r\Big]
\end{equation}
Dada \eqref{eq:v_ds minima} es posible cancelar:
    \[\frac{di_{ds}^r}{dt} = \frac{1}{L_d}\Big[\cancel{v_{ds}^r} - R_s(T_s^\circ(t))i_{ds}^r+ \cancel{L_qi_{ds}^r\omega_r}\Big]\]
Lo que resulta en una EDO de primer orden:
    \begin{equation}\label{eq:edo corriente id}
        \frac{di_{ds}^r(t)}{dt} = -\frac{1}{L_d} R_s(T_s^\circ(t))i_{ds}^r(t)
    \end{equation}
Resolviendo \eqref{eq:edo corriente id} se encuentra que \[ i_{ds}^r(t) = i_{ds}^r(0) \ e^{-\frac{R_s(t)}{L_d}\cdot t}\]
